\documentclass[11pt]{article}
\usepackage{amsmath, amssymb}

\title{CSE400 – Fundamentals of Probability in Computing}
\author{Scribe: Inequalities and Tail Bounds (Lectures 21–22)}
\date{Instructor: Dhaval Patel}

\begin{document}

\maketitle

\section*{1. Tails and Motivation}

\subsection*{Definition 1.1 (Tail of a Random Variable)}

The tail of a random variable $( X )$ is
\[
P\{X > x\}
\]

\subsection*{Motivation}

\begin{itemize}
\item Jobs delayed more than 24 hours
\item Packets dropped due to full buffer
\end{itemize}

\textbf{Key Idea:}

\begin{itemize}
\item Mean — average behavior (easy)
\item Tail — rare extreme events (hard)
\end{itemize}

\textbf{Why tails are hard to compute:}

\begin{itemize}
\item Requires summing many small probabilities
\end{itemize}

\section*{2. Tail Examples}

\subsection*{Example 2.1 (Binomial Tail)}

Suppose $( n )$ jobs are distributed among $( n )$ servers randomly.

\begin{itemize}
\item Let $( X \sim \text{Binomial}(n,p) )$
\item Probability of overload:
\[
P\{X > k\} = \sum \binom{n}{i} p^i (1-p)^{n-i}
\]
\end{itemize}

This is a tail event (rare overload).

\subsection*{Example 2.2 (Poisson Tail)}

Jobs arrive according to a Poisson process with rate $( \lambda )$.

\begin{itemize}
\item $( X \sim \text{Poisson}(\lambda) )$
\item Tail probability:
\[
P\{X > k\}
\]
\end{itemize}

This is a tail event (sudden spike in load).

\section*{3. Why Tail Bounds}

\subsection*{Observation}

Exact tail probabilities are often hard to compute:
\[
P\{X > k\} = \sum \cdots
\]

\textbf{Idea:}

\begin{itemize}
\item Instead of exact value, find an upper bound
\item Easier to compute
\item Still gives a guarantee
\end{itemize}

\subsection*{Running Example}

\begin{itemize}
\item $( X \sim \text{Binomial}(n, 1/2) )$
\item $( E[X] = n/2 )$
\item Question:
\[
P(X > 3n/4)
\]
\end{itemize}

\section*{4. Markov’s Inequality}

\subsection*{Key Idea}

\begin{itemize}
\item Uses only the mean
\item Large values $\Rightarrow$ rare
\end{itemize}

\subsection*{Theorem 4.1 (Markov’s Inequality)}

For a non-negative random variable $( X )$,
\[
P(X > a) < \frac{E[X]}{a}
\]

\subsection*{Proof (Idea)}

\begin{itemize}
\item Consider only $( X > a )$
\item Replace $( X )$ by $( a )$
\end{itemize}

\[
E[X] = \sum x \cdot p_X(x)
\]
\[
\ge \sum_{x \ge a} x \cdot p_X(x)
\]
\[
\ge a \cdot P(X > a)
\]
\[
\Rightarrow P(X > a) < \frac{E[X]}{a}
\]

\subsection*{Example 4.1 (Running Example)}

\begin{itemize}
\item $( X \sim \text{Binomial}(n,1/2) )$
\item $( E[X] = n/2 )$
\end{itemize}

\[
P(X > 3n/4) < \frac{E[X]}{3n/4}
\]

\subsection*{Solved Example 4.2 (Exam Grades)}

\begin{itemize}
\item $( E[X] = 70 )$, threshold $( a = 90 )$
\end{itemize}

Steps:

\begin{enumerate}
\item Let $( X > 0 )$, $( E[X] = 70 )$
\item Apply Markov:
\[
P(X > a) < \frac{E[X]}{a}
\]
\item Substitute:
\[
P(X > 90) < \frac{70}{90}
\]
\item Simplify:
\[
P(X > 90) < 0.778
\]
\end{enumerate}

\textbf{Result:} At most 77.8\% receive an A.

\subsection*{Solved Example 4.3 (Server Load)}

\begin{itemize}
\item $( E[X] = 5000 )$, $( a = 20000 )$
\end{itemize}

Steps:

\begin{enumerate}
\item $( X = )$ requests/sec
\item $( a = 20000 )$
\item Apply Markov:
\[
P(X > a) < \frac{E[X]}{a}
\]
\item Substitute:
\[
P(X > 20000) < \frac{5000}{20000}
\]
\end{enumerate}

\textbf{Result:}
\[
P(X > 20000) < 0.25
\]

\section*{5. Chebyshev’s Inequality}

\subsection*{Key Idea}

\begin{itemize}
\item Uses mean + variance
\item Controls deviation from mean
\end{itemize}

\subsection*{Theorem 5.1 (Chebyshev’s Inequality)}

\[
P(|X - \mu| > a) < \frac{\text{Var}(X)}{a^2}
\]

\subsection*{Proof (Idea)}

Apply Markov’s inequality to $( (X - \mu)^2 )$:
\[
P(|X - \mu| > a) = P((X-\mu)^2 > a^2)
\]
\[
< \frac{E[(X-\mu)^2]}{a^2} = \frac{\text{Var}(X)}{a^2}
\]

\subsection*{Example 5.1 (Running Example)}

\begin{itemize}
\item $( X \sim \text{Binomial}(n,1/2) )$
\item $( \mu = n/2 )$
\end{itemize}

\[
P(X > 3n/4) = P(|X - n/2| > n/4)
\]
\[
< \frac{\text{Var}(X)}{(n/4)^2}
\]

\subsection*{Example 5.2 (Exam Scores)}

\begin{itemize}
\item Mean = 70
\item Std deviation = 10
\end{itemize}

\[
P(|X - 70| > 30) < \frac{10^2}{30^2} = \frac{1}{9}
\]

\section*{6. Chernoff Bound}

\subsection*{Idea}

\begin{itemize}
\item Uses exponential transformation and Markov’s inequality
\end{itemize}

\subsection*{Lower Tail Derivation}

Steps:

\begin{enumerate}
\item Choose $( t < 0 )$:
\[
P(X < a) = P(tX > ta)
\]
\item Exponentiate:
\[
P(e^{tX} > e^{ta})
\]
\item Apply Markov:
\[
P(e^{tX} > e^{ta}) < \frac{E[e^{tX}]}{e^{ta}}
\]
\item Minimize over $( t < 0 )$
\end{enumerate}

\textbf{Result:}
\[
P(X < a) < \min_{t<0} e^{-ta} E[e^{tX}]
\]

\section*{7. Jensen’s Inequality}

\subsection*{Definition 7.1 (Convex Function)}

A function $( g(x) )$ is convex on interval $( S )$ if for any $( x_1, x_2 \in S )$ and $( \alpha \in [0,1] )$,
\[
g(\alpha x_1 + (1-\alpha)x_2) < \alpha g(x_1) + (1-\alpha)g(x_2)
\]

\subsection*{Convex Combination}

If $( \sum a_i = 1 )$,
\[
g\left(\sum a_i x_i\right) < \sum a_i g(x_i)
\]

\subsection*{Theorem 7.1 (Jensen’s Inequality)}

If $( g(x) )$ is convex and $( X )$ takes values in its domain,
\[
g(E[X]) < E[g(X)]
\]

\subsection*{Interpretation}

\begin{itemize}
\item Average first $\Rightarrow$ smaller value
\item Apply function first $\Rightarrow$ larger value
\item Convex functions push expectation upwards
\end{itemize}

\subsection*{Solved Example 7.1}

Let:

\begin{itemize}
\item $( X = 2 )$ with probability 0.5
\item $( X = 4 )$ with probability 0.5
\item $( g(x) = x^2 )$
\end{itemize}

Steps:

\begin{enumerate}
\item Compute expectation:
\[
E[X] = (2)(0.5) + (4)(0.5) = 3
\]
\item Left side:
\[
g(E[X]) = 3^2 = 9
\]
\item Distribution of $( g(X) )$:
\begin{itemize}
\item $( X^2 = 4 )$ with probability 0.5
\item $( X^2 = 16 )$ with probability 0.5
\end{itemize}
\item Right side:
\[
E[g(X)] = (4)(0.5) + (16)(0.5) = 10
\]
\end{enumerate}

\textbf{Result:}
\[
g(E[X]) = 9 < 10 = E[g(X)]
\]

\section*{8. System-Level Insight}

\subsection*{Comparison of Inequalities}

\begin{itemize}
\item Markov: uses only average
\item Chebyshev: uses average + variance
\item Chernoff: uses full distribution
\end{itemize}

\subsection*{Final Principle}

\textbf{More statistical information = Tighter bounds = Lower hardware costs}

\section*{End of Scribe}

\end{document}